\chapter{Integers}
\label{ch:integers}

In $\N$, $3 - 5$ doesn't exist. We build a larger system where it does.

The idea: every integer should be a difference $a - b$ of natural numbers, but we
can't write $a - b$ yet. So we use the pair $(a, b)$ as a stand-in for
\enquote{$a - b$}. Different pairs stand for the same difference: $(0, 2)$, $(1,
3)$ and $(5, 7)$ all mean \enquote{$-2$}. When do $(a, b)$ and $(c, d)$ mean the
same thing? When $a - b = c - d$, that is, when $a + d = c + b$. The last equation
uses only addition, which we have. So we \emph{define} an integer as a class of
pairs under that equation.

\section{Construction}

\begin{definition}
\label{def:integer-relation}
On $\N \times \N$ let
\[
  (a, b) \sim (c, d) \defiff a + d = c + b.
\]
\end{definition}

\begin{lemma}
\label{lem:integer-equivalence}
$\sim$ is an equivalence relation.
\end{lemma}

\begin{proof}
Reflexive: $a + b = a + b$. Symmetric: $a + d = c + b \Rightarrow c + b = a + d$.
Transitive: let $a + d = c + b$ and $c + f = e + d$. Adding the equations and
rearranging (\cref{thm:addition-laws}) gives $(a + f) + (c + d) = (e + b) + (c +
d)$. Cancelling $c + d$ (\cref{thm:addition-laws}(5)) gives $a + f = e + b$, so
$(a, b) \sim (e, f)$.
\end{proof}

\begin{definition}[Integers]
\label{def:integers}
$\Z \coloneqq (\N \times \N)/{\sim}$. Its elements are the \emph{integers}. The
class of $(a, b)$ is $[(a, b)]$; think of it as \enquote{$a - b$}.
\register{integers}{$\Z$}{the set of integers}
\register{integers}{$[(a, b)]$}{the integer \enquote{$a - b$}; replaced by $a - b$
after \cref{thm:integers-embedding}}
\end{definition}

By \cref{lem:classes}, $[(a, b)] = [(c, d)] \Leftrightarrow a + d = c + b$.

\section{Arithmetic}

The definitions below come from asking what $(a - b) + (c - d)$, $-(a - b)$ and
$(a - b) \cdot (c - d)$ \emph{would} be, and writing the answer as a difference
without using subtraction.

\begin{definition}[Operations on $\Z$]
\label{def:integer-operations}
\begin{align*}
  [(a, b)] + [(c, d)] &\coloneqq [(a + c,\; b + d)], \\
  -[(a, b)] &\coloneqq [(b, a)], \\
  [(a, b)] \cdot [(c, d)] &\coloneqq [(a \cdot c + b \cdot d,\; a \cdot d + b \cdot
  c)],
\end{align*}
with $0 \coloneqq [(0, 0)]$ and $1 \coloneqq [(1, 0)]$. \emph{Subtraction} is $x -
y \coloneqq x + (-y)$. The symbols $+$, $\cdot$, $0$ and $1$ are the same as in
$\N$; \cref{thm:integers-embedding} shows this is consistent
(\cref{rule:identification}).
\register{integers}{$x + y$, $x \cdot y$}{sum and product of integers (extend
those of $\N$)}
\register{integers}{$-x$}{the negative of $x$}
\register{integers}{$0$, $1$}{$[(0, 0)]$ and $[(1, 0)]$ (identified with
$0, 1 \in \N$)}
\end{definition}

These use representatives, so they are only valid definitions once they're shown
to be well-defined (\cref{def:well-defined,thm:map-rule}).

\begin{lemma}
\label{lem:integer-well-defined}
The right-hand sides in \cref{def:integer-operations} don't depend on the choice
of representatives.
\end{lemma}

\begin{proof}
Let $(a, b) \sim (a', b')$, so $a + b' = a' + b$, and $(c, d) \sim (c', d')$, so $c
+ d' = c' + d$.

\emph{Addition.} We need $(a + c) + (b' + d') = (a' + c') + (b + d)$. That is the
sum of the two hypotheses, rearranged.

\emph{Negation.} We need $b + a' = b' + a$. That is the first hypothesis, read
backwards.

\emph{Multiplication.} First change only the first factor. We need
\[
  (a \cdot c + b \cdot d) + (a' \cdot d + b' \cdot c) = (a' \cdot c + b' \cdot d) +
  (a \cdot d + b \cdot c).
\]
By \cref{thm:addition-laws,thm:multiplication-laws}, the left side is $(a + b')
\cdot c + (b + a') \cdot d$ and the right side is $(a' + b) \cdot c + (b' + a)
\cdot d$. These are equal because $a + b' = a' + b$. The formula is symmetric in
the two pairs, so the same argument covers changing the second factor. Changing
both is one change after the other.
\end{proof}

\begin{definition}[Commutative ring]
\label{def:ring}
A \emph{commutative ring} is a set $R$ with two binary operations $+$ and $\cdot$,
a unary operation $-$, and two elements $0, 1 \in R$ such that for all $x, y, z
\in R$:
\begin{center}
\begin{tabular}{@{}lll@{}}
  \toprule
  & Addition & Multiplication \\
  \midrule
  Associative & $(x + y) + z = x + (y + z)$ & $(x \cdot y) \cdot z = x \cdot (y \cdot z)$ \\
  Commutative & $x + y = y + x$ & $x \cdot y = y \cdot x$ \\
  Neutral element & $x + 0 = x$ & $x \cdot 1 = x$ \\
  Inverse & $x + (-x) = 0$ & (not required) \\
  Distributive & \multicolumn{2}{l}{$x \cdot (y + z) = x \cdot y + x \cdot z$} \\
  \bottomrule
\end{tabular}
\end{center}
Subtraction is $x - y \coloneqq x + (-y)$.
\register{rings}{$(R, +, \cdot, -, 0, 1)$}{a commutative ring}
\end{definition}

$\N$ is not a commutative ring, since additive inverses are missing.

\begin{theorem}
\label{thm:integers-ring}
$\Z$, with the operations of \cref{def:integer-operations}, is a commutative ring,
and $0 \neq 1$.
\end{theorem}

\begin{proof}
Write $x = [(a, b)]$, $y = [(c, d)]$ and $z = [(e, f)]$. Every law reduces to laws
of $\N$ (\cref{thm:addition-laws,thm:multiplication-laws}) applied to the
components.

\emph{Addition} is associative and commutative because it is componentwise. $x + 0
= [(a + 0, b + 0)] = x$. $x + (-x) = [(a + b, b + a)] = [(0, 0)]$, since $(a + b) +
0 = 0 + (b + a)$.

\emph{Multiplication.} Commutative: the formula is symmetric in the two pairs.
Neutral: $x \cdot 1 = [(a \cdot 1 + b \cdot 0,\; a \cdot 0 + b \cdot 1)] = [(a,
b)]$. Associative: expanding with distributivity in $\N$, $(x \cdot y) \cdot z$ and
$x \cdot (y \cdot z)$ both equal
\[
  [(a \cdot c \cdot e + a \cdot d \cdot f + b \cdot c \cdot f + b \cdot d \cdot
  e,\; a \cdot c \cdot f + a \cdot d \cdot e + b \cdot c \cdot e + b \cdot d \cdot
  f)].
\]

\emph{Distributivity.} Both $x \cdot (y + z)$ and $x \cdot y + x \cdot z$ equal
\[
  [(a \cdot c + a \cdot e + b \cdot d + b \cdot f,\; a \cdot d + a \cdot f + b
  \cdot c + b \cdot e)].
\]

$0 \neq 1$: $[(0, 0)] = [(1, 0)]$ would mean $0 + 0 = 1 + 0$, so $0 = 0 + 1$,
against \cref{thm:peano}(3).
\end{proof}

\begin{lemma}[Consequences of the ring laws]
\label{lem:ring-laws}
In every commutative ring:
\begin{enumerate}
  \item the additive inverse is unique: if $x + y = 0$, then $y = -x$;
  \item $x \cdot 0 = 0$;
  \item $(-x) \cdot y = -(x \cdot y)$ and $(-x) \cdot (-y) = x \cdot y$;
  \item $-(-x) = x$ and $-(x + y) = (-x) + (-y)$;
  \item $x + z = y + z \Rightarrow x = y$.
\end{enumerate}
\end{lemma}

\begin{proof}
(1) $y = y + (x + (-x)) = (y + x) + (-x) = 0 + (-x) = -x$.

(2) $x \cdot 0 = x \cdot (0 + 0) = x \cdot 0 + x \cdot 0$. Adding $-(x \cdot 0)$
to both sides gives $0 = x \cdot 0$.

(4) $(-x) + x = 0$, so $x = -(-x)$ by (1). And $(x + y) + ((-x) + (-y)) = 0$, so
(1) applies again.

(3) $(-x) \cdot y + x \cdot y = ((-x) + x) \cdot y = 0 \cdot y = 0$, so $(-x) \cdot
y = -(x \cdot y)$ by (1). Then $(-x) \cdot (-y) = -(x \cdot (-y)) = -(-(x \cdot
y)) = x \cdot y$ by (4).

(5) Add $-z$ to both sides.
\end{proof}

Since only the ring laws are used, \cref{lem:ring-laws} holds equally for $\Q$, and
later for $\R$, the complex numbers and rings of matrices. Proving things
\emph{once, from the laws}, is the main gain of naming a structure.

\section{\texorpdfstring{$\N$}{N} inside \texorpdfstring{$\Z$}{Z}}

\begin{theorem}[Embedding]
\label{thm:integers-embedding}
Let $\iota \colon \N \to \Z$ be the map with $n \via{\iota} [(n, 0)]$. Then $\iota$
is injective, $0 \via{\iota} 0$, $1 \via{\iota} 1$, and for all $m, n \in \N$ and
$x, y \in \Z$:
\[
  m \via{\iota} x \land n \via{\iota} y \;\Rightarrow\; m + n \via{\iota} x + y
  \;\land\; m \cdot n \via{\iota} x \cdot y \;\land\; [(m, n)] = x - y.
\]
\end{theorem}

\begin{proof}
Injective: $[(m, 0)] = [(n, 0)]$ means $m + 0 = n + 0$. Next, $[(m, 0)] + [(n, 0)]
= [(m + n, 0)]$ and $[(m, 0)] \cdot [(n, 0)] = [(m \cdot n + 0,\; m \cdot 0 + 0
\cdot n)] = [(m \cdot n, 0)]$. Finally $[(m, 0)] - [(n, 0)] = [(m, 0)] + [(0, n)]
= [(m, n)]$.
\end{proof}

So $\iota$ copies $\N$, with its arithmetic, into $\Z$. By
\cref{rule:identification} \emph{we identify every $n \in \N$ with the integer
$[(n, 0)]$} and treat $\N$ as a subset of $\Z$. Strictly speaking, the set $2 =
\set{0, 1}$ and the class $[(2, 0)]$ are different sets, but nothing we do depends
on that, because $\iota$ preserves all the structure. Under the identification
the construction says what we wanted: $[(a, b)] = a - b$. By
\cref{rule:abbreviation}, from now on integers are written with $-$ and the class
notation is not used again.

\begin{corollary}[Normal form]
\label{cor:integer-normal-form}
Every integer $x$ satisfies exactly one of
\[
  x = n \text{ for some } n \in \N \setminus \set{0}, \qquad x = 0, \qquad x = -n
  \text{ for some } n \in \N \setminus \set{0},
\]
and $n$ is unique.
\end{corollary}

\begin{proof}
Write $x = a - b$ with $a, b \in \N$. If $b \le a$, then $a = b + d$ for some $d \in
\N$ (\cref{thm:order-arithmetic}(1)), and $x = (b + d) - b = d$. If $a < b$, then
$b = a + d$ with $d \neq 0$, and $x = a - (a + d) = -d$. The cases are exclusive
and $n$ is unique: $n = m$ in $\Z$ gives $n = m$ in $\N$ because $\iota$ is
injective, and $n = -m$ gives $n + m = 0$, so $n = m = 0$
(\cref{thm:addition-laws}(6)).
\end{proof}

\section{No zero divisors}

\begin{theorem}
\label{thm:integral-domain}
For $x, y, z \in \Z$:
\begin{enumerate}
  \item $x \cdot y = 0 \Rightarrow x = 0 \lor y = 0$;
  \item $x \cdot z = y \cdot z \land z \neq 0 \Rightarrow x = y$ (cancellation).
\end{enumerate}
A commutative ring with $0 \neq 1$ and property (1) is called an \emph{integral
domain}.
\end{theorem}

\begin{proof}
(1) By \cref{cor:integer-normal-form}, $x$ is $m$ or $-m$ and $y$ is $n$ or $-n$
for some $m, n \in \N$. By \cref{lem:ring-laws}(3), $x \cdot y$ is $m \cdot n$ or
$-(m \cdot n)$. If $x \cdot y = 0$, then $m \cdot n = 0$, because $-k = 0$ only for
$k = 0$. So $m = 0$ or $n = 0$ (\cref{thm:multiplication-laws}(7)).

(2) $(x - y) \cdot z = x \cdot z - y \cdot z = 0$ by distributivity and
\cref{lem:ring-laws}(3). Since $z \neq 0$, (1) gives $x - y = 0$, so $x = y$.
\end{proof}

\section{Order}

\begin{definition}[Order on $\Z$]
\label{def:integer-order}
$x \le y \defiff y - x \in \N$. An integer $x$ is \emph{positive} if $0 < x$ and
\emph{negative} if $x < 0$.
\register{integers}{$x \le y$, $x < y$}{order of integers: $y - x \in \N$
(extends the order of $\N$)}
\end{definition}

\begin{theorem}[$\Z$ is an ordered ring]
\label{thm:integers-ordered}
\leavevmode
\begin{enumerate}
  \item $\le$ is a total order on $\Z$;
  \item $x \le y \Rightarrow x + z \le y + z$;
  \item $0 \le x \land 0 \le y \Rightarrow 0 \le x \cdot y$;
  \item on $\N$, this order agrees with the order of \cref{ch:naturals}.
\end{enumerate}
\end{theorem}

\begin{proof}
We use three facts: $\N$ is closed under $+$ and $\cdot$
(\cref{thm:integers-embedding}); for every $x$, $x \in \N$ or $-x \in \N$; and
$n \in \N$ together with $-n \in \N$ forces $n = 0$ (both by
\cref{cor:integer-normal-form}).

(1) Reflexive: $x - x = 0 \in \N$. Antisymmetric: if $y - x \in \N$ and $x - y =
-(y - x) \in \N$, then $y - x = 0$. Transitive: $z - x = (z - y) + (y - x)$. Total:
$y - x$ or $-(y - x) = x - y$ lies in $\N$.

(2) $(y + z) - (x + z) = y - x$.

(3) $x \cdot y \in \N$ when $x, y \in \N$.

(4) For $m, n \in \N$: $m \le n$ in $\N$ if and only if $n = m + d$ for some $d
\in \N$ (\cref{thm:order-arithmetic}(1)), if and only if $n - m \in \N$.
\end{proof}

Familiar rules follow from \cref{thm:integers-ordered}: for example $x \le y
\Leftrightarrow -y \le -x$, and $x \le y \land 0 \le z \Rightarrow x \cdot z \le y
\cdot z$ (Exercise~\ref{ex:integer-rules}).

\section{What is still missing}

In $\Z$ every equation $a + x = b$ has exactly one solution, $x = b - a$. But $2
\cdot x = 1$ has none: $2 \cdot x$ is $0$ for $x = 0$, at least $2$ for $0 < x$,
and negative for $x < 0$. \Cref{ch:rationals} adds the missing quotients.

\begin{exercises}
  \item Show that for $a, b, c, d \in \N$: $a - b \le c - d \Leftrightarrow a + d
    \le c + b$.
  \item Let $|\;|$ be the unary operation on $\Z$ with $|x| = x$ if $0 \le x$ and
    $|x| = -x$ otherwise. Prove $|x \cdot y| = |x| \cdot |y|$ and $|x + y| \le |x| +
    |y|$ (the \emph{triangle inequality}).
  \item \label{ex:integer-rules}
    Prove from \cref{thm:integers-ordered}: $x \le y \Rightarrow -y \le -x$; $x \le
    y \land 0 \le z \Rightarrow x \cdot z \le y \cdot z$; $0 \le x \cdot x$ for
    all $x$; $0 < 1$.
  \item Why does the construction need \emph{ordered} pairs? Try using unordered
    pairs $\set{a, b}$ instead of $(a, b)$ and find where it breaks. (Compare
    \enquote{$3 - 5$} and \enquote{$5 - 3$}.)
\end{exercises}
